\section{Fenwick 2D offline}
\whenuse{Altas/bajas de puntos y conteos o sumas dentro de rect\'angulos cuando
todas las coordenadas actualizables se conocen antes de procesar.}

Cada punto $(x,y)$ tiene un peso: el n\'umero que aporta a las sumas. Para
contar puntos, agregarlo usa peso $+1$ y eliminarlo peso $-1$; tambi\'en puede
representar dinero, poblaci\'on u otra cantidad asociada a la coordenada.

Se llama \emph{offline} porque primero se leen todas las operaciones para
conocer qu\'e coordenadas podr\'an actualizarse. Luego se reservan esos puntos,
se construye la estructura y se procesan las operaciones en su orden original.
No usa archivos. Cada nodo del Fenwick exterior sobre $x$ guarda otro Fenwick
con solamente las coordenadas $y$ que puede recibir.

\complexity{reserva y construcci\'on $O(n\log^2 n)$ en el peor caso; update y
consulta $O(\log^2 n)$; memoria $O(n\log n)$.}

\begin{lstlisting}[style=cpp]
struct OfflineFenwick2D {
    int size_x;
    vector<vector<int>> coordinates_y;
    vector<vector<ll>> bit;

    explicit OfflineFenwick2D(int x_count)
        : size_x(x_count), coordinates_y(x_count + 1), bit(x_count + 1) {}

    void reserve_point(int x, int y) {
        for (int index = x + 1; index <= size_x; index += index & -index)
            coordinates_y[index].push_back(y);
    }

    void build() {
        for (int index = 1; index <= size_x; ++index) {
            auto &values = coordinates_y[index];
            sort(values.begin(), values.end());
            values.erase(unique(values.begin(), values.end()), values.end());
            bit[index].assign(values.size() + 1, 0);
        }
    }

    void add(int x, int y, ll delta) {
        for (int index = x + 1; index <= size_x; index += index & -index) {
            const auto &values = coordinates_y[index];
            auto iterator = lower_bound(values.begin(), values.end(), y);
            assert(iterator != values.end() && *iterator == y);
            int position = iterator - values.begin() + 1;
            for (int inner = position; inner < (int)bit[index].size();
                 inner += inner & -inner)
                bit[index][inner] += delta;
        }
    }

    ll prefix_sum(int x_exclusive, int y_exclusive) const {
        ll answer = 0;
        for (int index = x_exclusive; index > 0; index -= index & -index) {
            int position = lower_bound(coordinates_y[index].begin(),
                                       coordinates_y[index].end(), y_exclusive)
                         - coordinates_y[index].begin();
            for (int inner = position; inner > 0; inner -= inner & -inner)
                answer += bit[index][inner];
        }
        return answer;
    }

    ll rectangle_sum(int x1, int y1, int x2, int y2) const {
        return prefix_sum(x2, y2) - prefix_sum(x1, y2)
             - prefix_sum(x2, y1) + prefix_sum(x1, y1);
    }
};
\end{lstlisting}

\subsection{Ejemplo de uso}

En un mapa aparecen y desaparecen objetos. Se pide contar cu\'antos hay dentro
de un rect\'angulo. Todas las coordenadas de futuras altas y bajas aparecen en
la entrada, por lo que pueden reservarse antes de procesarla.

\begin{lstlisting}[style=cpp]
OfflineFenwick2D points(5); // x validas: [0,5)

// Primera pasada: reservar todo punto que pueda actualizarse.
points.reserve_point(2, 5);
points.reserve_point(4, 3);
points.reserve_point(3, 6);
points.build();

// Segunda pasada: procesar las operaciones en orden.
points.add(2, 5, 1); // aparece un objeto
points.add(4, 3, 1);
points.add(3, 6, 1);

// Rectangulo [2,5) x [3,7): contiene los tres objetos.
cout << points.rectangle_sum(2, 3, 5, 7) << '\n'; // 3

points.add(4, 3, -1); // desaparece un objeto
cout << points.rectangle_sum(2, 3, 5, 7) << '\n'; // 2
\end{lstlisting}

\texttt{prefix\_sum(x,y)} suma los pesos de los puntos con coordenadas
$X<x$ e $Y<y$. Por inclusi\'on--exclusi\'on,
\texttt{rectangle\_sum(x1,y1,x2,y2)} consulta el rect\'angulo semiabierto
$[x_1,x_2)\times[y_1,y_2)$.

\warning{Las coordenadas $x$ deben estar comprimidas en $[0,size\_x)$. Llamar
\texttt{reserve\_point} para todo punto actualizable y despu\'es
\texttt{build} exactamente una vez. No intentar actualizar una coordenada que
no fue reservada.}
